% GENERATED FILE. Edit canonical lessons, then run npm run generate:books.
% canonical-source-hash: fnv1a64:caa24256a451a5d5
% canonical-lessons: JA-C178-butaniku, JA-C178-gyuuniku, JA-C178-toriniku, JA-C178-gyuunyuu, JA-C178-hiyasu

\chapter{Meat and Milk}
\label{ch:ja-meat-and-milk}
\hlchaptermodality{178}

\begin{chapteropening}
\textbf{By the end of this chapter:} \emph{I can name meats and milk and keep food cold.}

Complete the last lesson of chapter 178: I can name meats and milk and keep food cold.
\end{chapteropening}

\section[\texorpdfstring{\ja{ぶたにく} (butaniku)}{ぶたにく (butaniku)}]{\ja{ぶたにく} (butaniku) — pork}

\label{lesson:JA-C178-butaniku}



\begin{quote}
Before the new one: say the Japanese for to get burnt, then the Japanese for to go bad.
\end{quote}

\subsection*{You'll want to know: \ja{ぶたにく}}
\textbf{\ja{ぶたにく}} — \emph{butaniku} — \textquotedblleft{}pork\textquotedblright{}.

\begin{quote}
\textbf{\ja{にくやで} \ja{ぶたにくを} \ja{かいます。}} — \emph{nikuya de butaniku o kaimasu} — \textquotedblleft{}I buy pork at the butcher's.\textquotedblright{}
\end{quote}

\begin{cousinweb}
\textbf{\ja{ぶたにく}} (\emph{butaniku}) is built from native \emph{buta}, \textquotedblleft{}pig\textquotedblright{}, and Sino-Japanese \emph{niku}, \textquotedblleft{}meat\textquotedblright{}, a reading borrowed from Chinese (Middle Chinese \emph{nyuwk}) (JMdict, KANJIDIC2).
\end{cousinweb}

\begin{grammarlens}[title={What you've built}]
One more word for everyday life at A2.
\end{grammarlens}

\subsection*{Your turn}
\begin{itemize}
\raggedright
  \item \emph{Say it:} \emph{butaniku}
  \item \emph{Say it:} \emph{butaniku}, once more
  \item \emph{Say it:} \emph{nikuya de butaniku o kaimasu} — I buy pork at the butcher's.
  \item \emph{Recall:} say the Japanese for to get burnt, then the Japanese for to go bad, then say \emph{butaniku} again
\end{itemize}

\begin{tcolorbox}[breakable,colback=teal!4,colframe=teal!35!black,title={Before you move on}]
What does \ja{ぶたにく} mean? (\textquotedblleft{}Pork\textquotedblright{}.) Say it once more.
\end{tcolorbox}
\section[\texorpdfstring{\ja{ぎゅうにく} (gyūniku)}{ぎゅうにく (gyūniku)}]{\ja{ぎゅうにく} (gyūniku) — beef}

\label{lesson:JA-C178-gyuuniku}



\begin{quote}
Before the new one: say the Japanese for to go bad, then the Japanese for pork.
\end{quote}

\subsection*{You'll want to know: \ja{ぎゅうにく}}
\textbf{\ja{ぎゅうにく}} — \emph{gyūniku} — \textquotedblleft{}beef\textquotedblright{}.

\begin{quote}
\textbf{\ja{ぎゅうにくは} \ja{すこし} \ja{たかいです。}} — \emph{gyūniku wa sukoshi takai desu} — \textquotedblleft{}Beef is a little expensive.\textquotedblright{}
\end{quote}

\begin{cousinweb}
\textbf{\ja{ぎゅうにく}} (\emph{gyūniku}) is Sino-Japanese (kango): two Chinese characters read with their on'yomi, readings borrowed from Chinese. They are \emph{gyū}, \textquotedblleft{}cow\textquotedblright{} (Middle Chinese \emph{ngjuw}) and \emph{niku}, \textquotedblleft{}meat\textquotedblright{} (\emph{nyuwk}) (KANJIDIC2; Middle Chinese in Baxter's transcription).
\end{cousinweb}

\begin{grammarlens}[title={What you've built}]
One more word for everyday life at A2.
\end{grammarlens}

\subsection*{Your turn}
\begin{itemize}
\raggedright
  \item \emph{Say it:} \emph{gyūniku}
  \item \emph{Say it:} \emph{gyūniku}, once more
  \item \emph{Say it:} \emph{gyūniku wa sukoshi takai desu} — Beef is a little expensive.
  \item \emph{Recall:} say the Japanese for to go bad, then the Japanese for pork, then say \emph{gyūniku} again
\end{itemize}

\begin{tcolorbox}[breakable,colback=teal!4,colframe=teal!35!black,title={Before you move on}]
What does \ja{ぎゅうにく} mean? (\textquotedblleft{}Beef\textquotedblright{}.) Say it once more.
\end{tcolorbox}
\section[\texorpdfstring{\ja{とりにく} (toriniku)}{とりにく (toriniku)}]{\ja{とりにく} (toriniku) — chicken meat; poultry}

\label{lesson:JA-C178-toriniku}



\begin{quote}
Before the new one: say the Japanese for pork, then the Japanese for beef.
\end{quote}

\subsection*{You'll want to know: \ja{とりにく}}
\textbf{\ja{とりにく}} — \emph{toriniku} — \textquotedblleft{}chicken meat; poultry\textquotedblright{}.

\begin{quote}
\textbf{\ja{ばんごはんに} \ja{とりにくを} \ja{やきます。}} — \emph{bangohan ni toriniku o yakimasu} — \textquotedblleft{}I grill chicken for dinner.\textquotedblright{}
\end{quote}

\begin{cousinweb}
\textbf{\ja{とりにく}} (\emph{toriniku}) is built from native \emph{tori}, \textquotedblleft{}bird\textquotedblright{}, and Sino-Japanese \emph{niku}, \textquotedblleft{}meat\textquotedblright{}, a reading borrowed from Chinese (Middle Chinese \emph{nyuwk}) (JMdict, KANJIDIC2).
\end{cousinweb}

\begin{grammarlens}[title={What you've built}]
One more word for everyday life at A2.
\end{grammarlens}

\subsection*{Your turn}
\begin{itemize}
\raggedright
  \item \emph{Say it:} \emph{toriniku}
  \item \emph{Say it:} \emph{toriniku}, once more
  \item \emph{Say it:} \emph{bangohan ni toriniku o yakimasu} — I grill chicken for dinner.
  \item \emph{Recall:} say the Japanese for pork, then the Japanese for beef, then say \emph{toriniku} again
\end{itemize}

\begin{tcolorbox}[breakable,colback=teal!4,colframe=teal!35!black,title={Before you move on}]
What does \ja{とりにく} mean? (\textquotedblleft{}Chicken meat; poultry\textquotedblright{}.) Say it once more.
\end{tcolorbox}
\section[\texorpdfstring{\ja{ぎゅうにゅう} (gyūnyū)}{ぎゅうにゅう (gyūnyū)}]{\ja{ぎゅうにゅう} (gyūnyū) — milk; cow's milk}

\label{lesson:JA-C178-gyuunyuu}



\begin{quote}
Before the new one: say the Japanese for beef, then the Japanese for chicken meat.
\end{quote}

\subsection*{You'll want to know: \ja{ぎゅうにゅう}}
\textbf{\ja{ぎゅうにゅう}} — \emph{gyūnyū} — \textquotedblleft{}milk; cow's milk\textquotedblright{}.

\begin{quote}
\textbf{\ja{まいあさ} \ja{ぎゅうにゅうを} \ja{のみます。}} — \emph{maiasa gyūnyū o nomimasu} — \textquotedblleft{}I drink milk every morning.\textquotedblright{}
\end{quote}

\begin{cousinweb}
\textbf{\ja{ぎゅうにゅう}} (\emph{gyūnyū}) is Sino-Japanese (kango): two Chinese characters read with their on'yomi, readings borrowed from Chinese. They are \emph{gyū}, \textquotedblleft{}cow\textquotedblright{} (Middle Chinese \emph{ngjuw}) and \emph{nyū}, \textquotedblleft{}milk\textquotedblright{} (\emph{nyuX}) (KANJIDIC2; Middle Chinese in Baxter's transcription).
\end{cousinweb}

\begin{grammarlens}[title={What you've built}]
One more word for everyday life at A2.
\end{grammarlens}

\subsection*{Your turn}
\begin{itemize}
\raggedright
  \item \emph{Say it:} \emph{gyūnyū}
  \item \emph{Say it:} \emph{gyūnyū}, once more
  \item \emph{Say it:} \emph{maiasa gyūnyū o nomimasu} — I drink milk every morning.
  \item \emph{Recall:} say the Japanese for beef, then the Japanese for chicken meat, then say \emph{gyūnyū} again
\end{itemize}

\begin{tcolorbox}[breakable,colback=teal!4,colframe=teal!35!black,title={Before you move on}]
What does \ja{ぎゅうにゅう} mean? (\textquotedblleft{}Milk; cow's milk\textquotedblright{}.) Say it once more.
\end{tcolorbox}
\section[\texorpdfstring{\ja{ひやす} (hiyasu)}{ひやす (hiyasu)}]{\ja{ひやす} (hiyasu) — to chill; to cool}

\label{lesson:JA-C178-hiyasu}



\begin{quote}
Before the new one: say the Japanese for chicken meat, then the Japanese for milk.
\end{quote}

\subsection*{You'll want to know: \ja{ひやす}}
\textbf{\ja{ひやす}} — \emph{hiyasu} — \textquotedblleft{}to chill; to cool\textquotedblright{}.

\begin{quote}
\textbf{\ja{なつは} \ja{すいかを} \ja{ひやします。}} — \emph{natsu wa suika o hiyashimasu} — \textquotedblleft{}In summer we chill watermelon.\textquotedblright{}
\end{quote}

Polite \textbf{\ja{ひやします}}, \emph{hiyashimasu}; a godan verb.

\begin{cousinweb}
\textbf{\ja{ひやす}} (\emph{hiyasu}) is a native Japanese word (wago), inherited rather than borrowed from Chinese; its kanji was attached later as a kun'yomi (JMdict, KANJIDIC2).
\end{cousinweb}

\begin{grammarlens}[title={What you've built}]
One more word for everyday life at A2.
\end{grammarlens}

\subsection*{Your turn}
\begin{itemize}
\raggedright
  \item \emph{Say it:} \emph{hiyasu}
  \item \emph{Say it:} \emph{hiyasu}, once more
  \item \emph{Say it:} \emph{natsu wa suika o hiyashimasu} — In summer we chill watermelon.
  \item \emph{Recall:} say the Japanese for chicken meat, then the Japanese for milk, then say \emph{hiyasu} again
\end{itemize}

\begin{tcolorbox}[breakable,colback=teal!4,colframe=teal!35!black,title={Before you move on}]
What does \ja{ひやす} mean? (\textquotedblleft{}To chill; to cool\textquotedblright{}.) Say it once more.
\end{tcolorbox}
